% Generated draft: requires mathematical and source review.
\documentclass[11pt]{article}
\usepackage[margin=25mm]{geometry}
\usepackage{fontspec}
\setmainfont{Latin Modern Roman}
\usepackage{amsmath,amssymb,mathtools,unicode-math}
\setmathfont{Latin Modern Math}
\usepackage{enumitem,needspace,xurl,hyperref}
\setlength{\parindent}{0pt}
\setlength{\parskip}{5pt}
\setlength{\emergencystretch}{3em}
\newcommand{\N}{\mathbb{N}}
\newcommand{\Z}{\mathbb{Z}}
\newcommand{\Q}{\mathbb{Q}}
\newcommand{\R}{\mathbb{R}}
\newcommand{\C}{\mathbb{C}}
\newcommand{\F}{\mathbb{F}}
\newcommand{\A}{\mathbb{A}}
\newcommand{\PP}{\mathbb{P}}
\newcommand{\E}{\mathbb{E}}
\newcommand{\eps}{\varepsilon}
\newcommand{\dd}{\,\mathrm{d}}
\begin{document}
\section*{UnsolvedMath ID: 701634}
\textbf{Source problem code:} catalogue entry
\subsection{Definitions and mathematical statement}
The supplied record uses notation whose definitions are referred to Problem 19.84 in the Archive, but that entry is not supplied. The following formalization uses the conventional meanings stated explicitly below; their agreement with the referenced entry requires verification.

Let \(\mathbb P\) be the set of primes, and let \(\sigma=\{\sigma_i:i\in I\}\) be a partition of \(\mathbb P\) into nonempty, pairwise disjoint subsets. For a positive integer \(n\), write \(\pi(n)\) for its set of prime divisors, with \(\pi(1)=\varnothing\). For a finite group \(X\), put \(\pi(X)=\pi(|X|)\). A subgroup \(H\leq X\) is a \(\sigma_i\)-Hall subgroup if
\[
\pi(H)\subseteq\sigma_i
\quad\text{and}\quad
\pi([X:H])\cap\sigma_i=\varnothing.
\]
A finite group \(X\) is \(\sigma\)-primary if \(\pi(X)\subseteq\sigma_j\) for some \(j\in I\). For \(B\leq X\), its core is
\[
\operatorname{Core}_X(B)=\bigcap_{x\in X}x^{-1}Bx.
\]
A subgroup \(A\leq G\) is \(\sigma\)-subnormal in \(G\) if there is a finite chain
\[
A=A_0\leq A_1\leq\cdots\leq A_n=G
\]
with \(n\geq0\), such that for every \(1\leq k\leq n\), either \(A_{k-1}\trianglelefteq A_k\) or
\[
A_k/\operatorname{Core}_{A_k}(A_{k-1})
\quad\text{is }\sigma\text{-primary}.
\]

Under these conventions, the question is whether the following implication holds for every such partition \(\sigma\), every finite group \(G\), and every subgroup \(A\leq G\): if \(G\) has a \(\sigma_i\)-Hall subgroup for every \(i\in I\), and
\[
\text{for every }i\in I\text{ and every }\sigma_i\text{-Hall subgroup }H\leq G,
\quad A\cap H\text{ is a }\sigma_i\text{-Hall subgroup of }A,
\]
then \(A\) is \(\sigma\)-subnormal in \(G\).

The supplied source reports an affirmative answer in each of two cases: when \(\sigma=\{\{p\}:p\in\mathbb P\}\), and when \(G\) is \(\sigma\)-soluble. Here the conventional meaning of \(\sigma\)-soluble is that every chief factor of \(G\) is \(\sigma\)-primary. A chief factor is a quotient \(N/M\), where \(M\trianglelefteq G\), \(N\trianglelefteq G\), \(M<N\), and there is no normal subgroup of \(G\) strictly between \(M\) and \(N\). These affirmative cases are supplied source qualifications, not independently verified results.
\subsection{Short English statement}
For a finite group having Hall subgroups for every block of a prime partition \(\sigma\), suppose a subgroup \(A\) intersects every such Hall subgroup in the corresponding Hall subgroup of \(A\). Must \(A\) be \(\sigma\)-subnormal? The source reports affirmative answers for the partition into singleton primes and for \(\sigma\)-soluble groups. The terminology used here requires confirmation against the unsupplied Archive Problem 19.84.
\subsection{Sources}
\begin{itemize}
\item Supplied local UnsolvedMath registry record: permanent ID \(701634\), exact source problem code \texttt{catalogue entry}, title \emph{Kourovka Notebook Problem 19.86}. The registry labels the problem open; this is supplied status metadata, not an independently verified assessment.
\item \emph{The Kourovka Notebook (Unsolved Problems in Group Theory)}, Problem 19.86, proposed by A. N. Skiba, 19th issue (2018). The supplied background also identifies the 21st issue edition (2026), arXiv:1401.0300v48. The supplied collection description names E. I. Khukhro and V. D. Mazurov as editors.
\item The supplied statement refers to Problem 19.84 in the Archive for terminology or context. That entry is not included in the supplied evidence. The statement also reports the two affirmative special cases, without supplying separate references for them.
\end{itemize}
\end{document}
